% Options for packages loaded elsewhere
\PassOptionsToPackage{unicode}{hyperref}
\PassOptionsToPackage{hyphens}{url}
\PassOptionsToPackage{dvipsnames,svgnames,x11names}{xcolor}
\documentclass[
  a4paper,
]{article}
\usepackage{xcolor}
\usepackage[margin=25mm]{geometry}
\usepackage{amsmath,amssymb}
\setcounter{secnumdepth}{-\maxdimen} % remove section numbering
\usepackage{iftex}
\ifPDFTeX
  \usepackage[T1]{fontenc}
  \usepackage[utf8]{inputenc}
  \usepackage{textcomp} % provide euro and other symbols
\else % if luatex or xetex
  \usepackage{unicode-math} % this also loads fontspec
  \defaultfontfeatures{Scale=MatchLowercase}
  \defaultfontfeatures[\rmfamily]{Ligatures=TeX,Scale=1}
\fi
\usepackage{lmodern}
\ifPDFTeX\else
  % xetex/luatex font selection
  \setmainfont[]{DejaVu Serif}
  \setsansfont[]{DejaVu Sans}
  \setmonofont[]{DejaVu Sans Mono}
\fi
% Use upquote if available, for straight quotes in verbatim environments
\IfFileExists{upquote.sty}{\usepackage{upquote}}{}
\IfFileExists{microtype.sty}{% use microtype if available
  \usepackage[]{microtype}
  \UseMicrotypeSet[protrusion]{basicmath} % disable protrusion for tt fonts
}{}
\makeatletter
\@ifundefined{KOMAClassName}{% if non-KOMA class
  \IfFileExists{parskip.sty}{%
    \usepackage{parskip}
  }{% else
    \setlength{\parindent}{0pt}
    \setlength{\parskip}{6pt plus 2pt minus 1pt}}
}{% if KOMA class
  \KOMAoptions{parskip=half}}
\makeatother
\usepackage{longtable,booktabs,array}
\newcounter{none} % for unnumbered tables
\usepackage{calc} % for calculating minipage widths
% Correct order of tables after \paragraph or \subparagraph
\usepackage{etoolbox}
\makeatletter
\patchcmd\longtable{\par}{\if@noskipsec\mbox{}\fi\par}{}{}
\makeatother
% Allow footnotes in longtable head/foot
\IfFileExists{footnotehyper.sty}{\usepackage{footnotehyper}}{\usepackage{footnote}}
\makesavenoteenv{longtable}
\usepackage{graphicx}
\makeatletter
\newsavebox\pandoc@box
\newcommand*\pandocbounded[1]{% scales image to fit in text height/width
  \sbox\pandoc@box{#1}%
  \Gscale@div\@tempa{\textheight}{\dimexpr\ht\pandoc@box+\dp\pandoc@box\relax}%
  \Gscale@div\@tempb{\linewidth}{\wd\pandoc@box}%
  \ifdim\@tempb\p@<\@tempa\p@\let\@tempa\@tempb\fi% select the smaller of both
  \ifdim\@tempa\p@<\p@\scalebox{\@tempa}{\usebox\pandoc@box}%
  \else\usebox{\pandoc@box}%
  \fi%
}
% Set default figure placement to htbp
\def\fps@figure{htbp}
\makeatother
\setlength{\emergencystretch}{3em} % prevent overfull lines
\providecommand{\tightlist}{%
  \setlength{\itemsep}{0pt}\setlength{\parskip}{0pt}}
\usepackage{luatexja}
\usepackage{xurl}
\usepackage{bookmark}
\IfFileExists{xurl.sty}{\usepackage{xurl}}{} % add URL line breaks if available
\urlstyle{same}
\hypersetup{
  pdftitle={Relational Readout of Quaternions and Crystallization as Inflation},
  pdfauthor={Noriaki Kihara (WF System Co., Ltd.) ORCID 0009-0004-6753-4020},
  colorlinks=true,
  linkcolor={Maroon},
  filecolor={Maroon},
  citecolor={Blue},
  urlcolor={Blue},
  pdfcreator={LaTeX via pandoc}}

\title{Relational Readout of Quaternions and Crystallization as
Inflation}
\usepackage{etoolbox}
\makeatletter
\providecommand{\subtitle}[1]{% add subtitle to \maketitle
  \apptocmd{\@title}{\par {\large #1 \par}}{}{}
}
\makeatother
\subtitle{― Anonymity-Preserving Autonomous Structuring and Partial
Projection ―}
\author{Noriaki Kihara (WF System Co., Ltd.) ORCID 0009-0004-6753-4020}
\date{October 21, 2026 Version DOI 10.5281/zenodo.23523238}

\ltjsetparameter{jacharrange={-2,-3}}
\begin{document}
\maketitle

\textbf{Author:} Noriaki Kihara (WF System Co., Ltd.)　ORCID
0009-0004-6753-4020\\
\textbf{Date:} 2026-10-21　　\textbf{Version:} v1.1\\
\textbf{Version DOI:} 10.5281/zenodo.23523238　　\textbf{Concept DOI:}
10.5281/zenodo.23523237\\
\textbf{Same folder:} figures fig01--fig09 (SVG; English-label versions
in \texttt{図表/\allowbreak{}en/\allowbreak{}}), figure-generation programs
\texttt{make\_\allowbreak{}figures\_\allowbreak{}quaternion\_\allowbreak{}crystallization\_\allowbreak{}20261021.py} and
\texttt{make\_\allowbreak{}figures\_\allowbreak{}quaternion\_\allowbreak{}crystallization\_\allowbreak{}20261021\_\allowbreak{}v1.1.py}\\
\textbf{Note:} The Japanese version
(\texttt{quaternion\_\allowbreak{}crystallization\_\allowbreak{}inflation\_\allowbreak{}ja\_\allowbreak{}v1.1.md}) is the
authoritative text; this English version is a translation.

\begin{center}\rule{0.5\linewidth}{0.5pt}\end{center}

\subsection{Abstract}\label{abstract}

The author's research programme aims to read out spacetime and particles
only from the relations within an anonymous set of complex numbers
\{zₙ\}, without placing a background spacetime, a ruler, or particle
species from outside, and to connect this to the fact that general
relativity and the Standard Model hold locally with extremely high
accuracy. The previous paper {[}S6{]} combined two complex numbers into
a quaternion q = z₁ + z₂ j and, since its polar form is exactly the
central projection onto the tangent hyperplane at the north pole,
constructed the map \{zₙ\} → \{qₙ\} → (1+3) as a partial projection.
However, an anonymous set of quaternions has no structure by itself, and
a mechanism by which structure arises autonomously, without supplying
from outside which quaternion relates to which, remained open.

This paper proposes the hypothesis that this mechanism is
crystallization in the sense of condensed-matter physics. The main
contents are as follows.

\begin{enumerate}
\def\labelenumi{\arabic{enumi}.}
\tightlist
\item
  From the relation between two states, a readout that takes the visual
  angle as distance is defined. For two equal-amplitude complex numbers
  we obtain the coupling coefficient g = \textbar z₁ − z₂\textbar{} /
  \textbar z₁ + z₂\textbar{} = tan(Δθ/2) and the angular-size distance D
  = 1/g = cot(Δθ/2). In phase is at infinity (no action); antiphase is
  contact (total reflection).
\item
  For quaternions, the coupling coefficient k = (q₁ + q₂)⁻¹(q₁ − q₂) =
  −n tan(α/2) becomes a three-dimensional vector. k is the stereographic
  projection of the relative phase u = q₁⁻¹q₂ as seen from the observer,
  and at the same time the central projection of the square root of the
  relative phase e\^{}(nα/2). A relation with coupling strength 1
  (orthogonality) corresponds, in the language of the previous paper, to
  a relation on the light cone.
\item
  Since this readout preserves angles, a localized observer cannot
  distinguish uniform curvature, global scale, the choice of map, or the
  choice of left/right frame. Local flatness is not an approximation but
  a property of the readout. What is observable is only
  direction-dependent distortion (tides).
\item
  As a way of not supplying connectivity from outside, the convex hull
  of an anonymous point set is used. Unit quaternions are points on S³ ⊂
  ℝ⁴, so their convex hull gives a truss of tetrahedra, which becomes a
  Delaunay triangulation under observer-centred stereographic
  projection. The test is invariant under conformal maps.
\item
  The non-uniqueness of the triangulation is taken not as a defect but
  as a principle. Under the rigidity count F − S = 3V − 6 − E, only
  partial projection, between zero projection (no structure forms) and
  full projection (rigid crystal, no evolution), combines a skeleton
  with room for rearrangement.
\item
  Since the structure is not unique, the composition assumption itself
  requires a process (evolution). When a seed enters a supercooled,
  highly symmetric state, the number of bondable faces multiplies by r
  (r ≤ 3) per process, and geometric crystallization occurs. For the
  length scale a ∝ V\^{}(1/3), the number of e-folds is N\_e = (n/3) ln
  r; for r = 3, the number of processes needed for 60 e-folds is n ≈
  164. The end is given structurally as faces close and r approaches 1.
\item
  Regular tetrahedra cannot fill flat space (gap 7.36°), but in S³ the
  600-cell fills it completely, and its 120 vertices are unit
  quaternions. In a locally flat readout, the part that cannot close
  necessarily remains as defect lines, and partial projection is
  enforced without external parameters. The deficit angles of the defect
  lines carry curvature.
\item
  The author's spacetime equation x² + y² + z² + t² = R² + Q₁² + Q₂² +
  Q₃² is the condition that the norm of a split octonion vanishes, and
  is equivalent to the two quaternions of the spacetime side and the
  internal side having equal amplitude.
\end{enumerate}

This is a hypothesis paper; the verification plan and falsification
conditions are given in §13.

\textbf{Keywords:} quaternions, anonymity, partial projection,
stereographic projection, Delaunay triangulation, rigidity, geometrical
frustration, 600-cell, supercooling, inflation, split octonions

\begin{center}\rule{0.5\linewidth}{0.5pt}\end{center}

\subsection{1. Research motivation}\label{research-motivation}

\subsubsection{1.1 The starting question}\label{the-starting-question}

General relativity and the Standard Model cannot be unified well, yet
when a locally flat approximation is made they hold with extremely high
accuracy. Why? Is there a way to explain this without introducing many
arbitrary parameters? This is the starting point of this research.

\subsubsection{1.2 History: each impasse forced the next
stage}\label{history-each-impasse-forced-the-next-stage}

The research proceeded in such a way that each impasse forced the next
stage (Figure 1).

\begin{figure}
\centering
\pandocbounded{\includegraphics[keepaspectratio,alt={History of the research. Red text marks the impasse at each stage, which forced the next stage.}]{figs/fig01_research_history_en.png}}
\caption{History of the research. Red text marks the impasse at each
stage, which forced the next stage.}
\end{figure}

\begin{itemize}
\tightlist
\item
  \textbf{Stage ①} Drawing on experience in CG and 3D modelling, the
  author studied the nonlinear map of central (gnomonic) projection and
  noticed that, with the map Σxₙ² = R², the curvature radius R of the
  central projection may be connected to the light cone {[}S1{]}.
  However, Σxₙ² = 0 has only the trivial solution over the reals.
\item
  \textbf{Stage ②} Therefore xₙ was taken to be complex, zₙ, and Σzₙ² =
  0 was imposed. The foundations of theoretical physics are built on
  complex Hilbert spaces and their mapping operations, but how to read
  them out as four-dimensional spacetime is not clear, and a background
  four-dimensional spacetime is assumed from the start. Also, a pure
  complex number can be read directly as a wave in exponential form, but
  as it stands it cannot represent a travelling wave, and the phase ωt
  is brought in implicitly. A question remained: where in a complex wave
  with only two degrees of freedom does the capacity to store
  information lie?
\item
  \textbf{Stage ③} From research on the light cone and the conjugacy of
  λν, the author noticed that spacetime can be written as x² + y² + z² −
  t² = R² + Q², that putting t as it gives x² + y² + z² + t² = R² + Q²,
  and that the Lorentz transformation and the Wick rotation are problems
  of reading out t; the author now considers it likely that x² + y² + z²
  + t² = R² + Q₁² + Q₂² + Q₃². At the same time, the author realized
  that anonymity is extremely important. The zₙ in Σzₙ² = 0 are
  anonymous, and the labels x, y, z, t, R, Q₁, Q₂, Q₃ are names attached
  afterwards to quantities that can be read as such when read out {[}S2,
  S3{]}. However, the set of complex numbers \{zₙ\} itself has no
  structure.
\item
  \textbf{Stage ④} The complex numbers were considered as a complete
  network, and a method of reading out structure from the M = N(N−1)/2
  relations for N vertices was examined {[}S4{]}. However, which vertex
  relates to which must be supplied from outside, which violates
  anonymity.
\item
  \textbf{Stage ⑤} Considering that the M waves cause M² actions
  including self-interaction, the author found that repeating this
  self-interaction in a set of complex numbers with very high symmetry
  produces inflation-like geometric growth {[}S5{]}. However,
  four-dimensional spacetime could not be read out of that process.
\item
  \textbf{Stage ⑥} Two complex numbers make a quaternion, and four make
  an octonion (counting from the reals, 1 → 2 → 4 → 8). A complex number
  is originally just a two-dimensional vector, and the names real part
  and imaginary part were also attached afterwards; it looked as if it
  would violate anonymity but did not. Likewise, a quaternion is just a
  four-dimensional vector and an octonion an eight-dimensional vector;
  they have structure (a product) but do not violate anonymity {[}S6{]}.
  Why physics chooses octonions is not asked in this paper; it is placed
  as an axiom.
\item
  \textbf{Stage ⑦ (this paper)} Can quaternions be used to elucidate a
  mechanism by which structure arises autonomously while preserving
  anonymity and without supplying relations from outside? The clue is
  crystallization in condensed-matter physics. When a slight seed is
  given to a supercooled structure of high symmetry, a crystal or a
  glass is constructed geometrically.
\end{itemize}

Quaternions and crystallization are not a whim; they were selected by
the constraint ``obtain structure while preserving anonymity.''

\subsubsection{1.3 Structure of this
paper}\label{structure-of-this-paper}

§2 sets the conventions. §3--§4 read out distance from the relation
between two states (complex numbers, quaternions), and §5 discusses the
localized observer and local flatness. §6 is a hypothesis about standing
waves and mass. §7 reads out connectivity from an anonymous point set,
and §8 expresses partial projection by the rigidity count. §9 discusses
frustration and defects, and §10 the correspondence between geometric
growth and inflation. §11 is a bridge to octonions, §12 compares with
previous research, and §13 gives the verification plan and falsification
conditions.

\begin{center}\rule{0.5\linewidth}{0.5pt}\end{center}

\subsection{2. Conventions}\label{conventions}

\subsubsection{2.1 What is inherited from the previous
paper}\label{what-is-inherited-from-the-previous-paper}

The axioms and definitions of the previous paper {[}S6{]} and the
catalogue {[}S3{]} are used as they are: anonymity (elements carry no
labels, and there is no special origin, coordinate, distance, sign, or
species), self-consistency (the system takes itself as the reference),
indistinguishability (a readout cannot distinguish which layer a value
belongs to), the linear closure Σq = 0 and the true-square closure Σq² =
0. The notation for quaternions (q = s + v, polar form q =
\textbar q\textbar{} e\^{}(nθ), conjugate q̄, true square q² = (s² −
\textbar v\textbar²) + 2s v) also follows the previous paper.

\subsubsection{2.2 Declaration of division, origin and
observer}\label{declaration-of-division-origin-and-observer}

This system is a system of products. Division is an implicit
normalization, so it is performed only once at the end of the readout,
and each time the denominator is declared. Addition inside products
moves the origin, so each time addition is used the origin is declared.
Every value carries the observer whose value it is. For the main
quantities of this paper, the declarations are as follows.

{\def\LTcaptype{none} % do not increment counter
\begin{longtable}[]{@{}
  >{\raggedright\arraybackslash}p{(\linewidth - 6\tabcolsep) * \real{0.2500}}
  >{\raggedright\arraybackslash}p{(\linewidth - 6\tabcolsep) * \real{0.2500}}
  >{\raggedright\arraybackslash}p{(\linewidth - 6\tabcolsep) * \real{0.2500}}
  >{\raggedright\arraybackslash}p{(\linewidth - 6\tabcolsep) * \real{0.2500}}@{}}
\toprule\noalign{}
\begin{minipage}[b]{\linewidth}\raggedright
Quantity
\end{minipage} & \begin{minipage}[b]{\linewidth}\raggedright
Denominator
\end{minipage} & \begin{minipage}[b]{\linewidth}\raggedright
Origin
\end{minipage} & \begin{minipage}[b]{\linewidth}\raggedright
Observer
\end{minipage} \\
\midrule\noalign{}
\endhead
\bottomrule\noalign{}
\endlastfoot
Relative phase u = q₁⁻¹q₂ & q₁ (left division) & --- & q₁ \\
Coupling coefficient k = (q₁ + q₂)⁻¹(q₁ − q₂) & the composite q₁ + q₂ &
the origin O of the sum and difference (the barycentre fixed by Σq = 0)
& q₁ \\
Distance x = k⁻¹ & k & same as above & q₁ \\
Amplitude ratio ρ = \textbar q₂\textbar{} / \textbar q₁\textbar{} &
\textbar q₁\textbar{} & --- & q₁ \\
Relative phase of a split octonion u = a⁻¹b (§11) & the spacetime side a
(left division) & the common origin O of a and b & a \\
\end{longtable}
}

Whether one reads with left division or right division is a choice that
enters from outside. This paper states it explicitly as the choice of
left or right, and describes how the results depend on that choice.

\begin{center}\rule{0.5\linewidth}{0.5pt}\end{center}

\subsection{3. Reading out distance from the relation between two
states: complex
numbers}\label{reading-out-distance-from-the-relation-between-two-states-complex-numbers}

\subsubsection{3.1 Visual angle and the isosceles
triangle}\label{visual-angle-and-the-isosceles-triangle}

When two complex numbers z₁, z₂ have equal amplitude
\textbar z₁\textbar{} = \textbar z₂\textbar{} = r and phase difference
Δθ, the origin O and z₁, z₂ form an isosceles triangle. Its shape is
fixed by the two diagonals (Figure 2).

\begin{itemize}
\tightlist
\item
  Difference: \textbar z₁ − z₂\textbar{} = 2r sin(Δθ/2)
\item
  Sum: \textbar z₁ + z₂\textbar{} = 2r cos(Δθ/2)
\end{itemize}

When something of size A is seen at visual angle Δθ, the height of the
isosceles triangle with it as base and the observer at the apex is the
distance D\_A = (A/2) cot(Δθ/2), in the same sense as the
angular-diameter distance of astronomy. The smaller the visual angle,
the farther; the larger the size, the farther.

\begin{figure}
\centering
\pandocbounded{\includegraphics[keepaspectratio,alt={The isosceles triangle formed by two equal-amplitude complex numbers (left), and the coupling coefficient g and the angular-size distance D (right).}]{figs/fig02_angular_distance_coupling_en.png}}
\caption{The isosceles triangle formed by two equal-amplitude complex
numbers (left), and the coupling coefficient g and the angular-size
distance D (right).}
\end{figure}

\subsubsection{3.2 Coupling coefficient and angular-size
distance}\label{coupling-coefficient-and-angular-size-distance}

From the ratio of the sum and the difference we define the dimensionless
quantities

g = \textbar z₁ − z₂\textbar{} / \textbar z₁ + z₂\textbar{} =
tan(Δθ/2),　D = 1/g = cot(Δθ/2).

D is the angular-size distance in units of the size A. g and D do not
change when both are multiplied by the same complex number (common
magnification and common phase rotation). The global scale and the
origin of phase cannot be read, so these quantities extract only the
relation.

\begin{itemize}
\tightlist
\item
  In phase (Δθ → 0): g → 0, D → ∞. The two appear superposed, have no
  action, and are at infinity.
\item
  Orthogonal (Δθ = π/2): g = D = 1. The sum and the difference balance.
\item
  Antiphase (Δθ = π): g → ∞, D = 0. With equal amplitude z₁ + z₂ = 0,
  and the cancellation corresponds exactly to contact.
\end{itemize}

If distance is read as the weakness of the coupling of action, this
behaviour is natural. Taking k = (z₁ − z₂)/(z₁ + z₂) as a complex
number, it has the same form as the reflection coefficient of a
transmission line Γ = (Z − Z₀)/(Z + Z₀): identical things are matched
with no reflection (no action), and antiphase is total reflection
(maximal action). It may also be read as the ratio of the differential
component (z₁ − z₂)/2 to the common-mode component (z₁ + z₂)/2.

\subsubsection{3.3 Two maps}\label{two-maps}

The relative phase u = z₂/z₁ = e\^{}(iΔθ) is a point on the unit circle,
and k = (1 − u)/(1 + u) = −i tan(Δθ/2). This is the coordinate obtained
by projecting the unit circle from the point u = −1 onto the imaginary
axis (the Cayley transform, stereographic projection), and the observer
itself, u = +1, comes to the origin k = 0 (Figure 3).

Conversely, the coordinate projected from the point u = +1, x = (1 +
u)/(1 − u) = 1/k, is the map of distances with magnitude cot(Δθ/2) = D.
The observer itself is sent to infinity. The self-distance is infinite,
and one's self is not a reference for distance.

The two maps are related by the inversion x = 1/k, and both preserve
angles.

\begin{figure}
\centering
\pandocbounded{\includegraphics[keepaspectratio,alt={The circle of relative phase and the two maps. Projecting from −1 gives the coupling map k; projecting from +1 gives the distance map x.}]{figs/fig03_relative_phase_two_maps_en.png}}
\caption{The circle of relative phase and the two maps. Projecting from
−1 gives the coupling map k; projecting from +1 gives the distance map
x.}
\end{figure}

\subsubsection{3.4 Composition law}\label{composition-law}

Chaining relative phases in the same plane, t = tan(Δθ/2) composes as

t₁₂ = (t₁ + t₂) / (1 − t₁t₂).

This is of the same type as the relativistic velocity addition and the
cascade connection of reflection coefficients. What is additive is the
phase difference Δθ; the coupling coefficient is not a quantity to be
added.

\begin{center}\rule{0.5\linewidth}{0.5pt}\end{center}

\subsection{4. Extension to quaternions}\label{extension-to-quaternions}

\subsubsection{4.1 The coupling coefficient becomes a three-dimensional
vector}\label{the-coupling-coefficient-becomes-a-three-dimensional-vector}

For two quaternions q₁, q₂, let the relative phase be u = q₁⁻¹q₂ and
define the coupling coefficient

k = (q₁ + q₂)⁻¹(q₁ − q₂) = (1 + u)⁻¹(1 − u).

Since q₁ + q₂ = q₁(1 + u) and q₁ − q₂ = q₁(1 − u), and 1 and u commute,
this expression is determined unambiguously as a function of u.

With equal amplitude, u = e\^{}(nα) (n a unit pure imaginary
quaternion), and the computation with the complex i replaced by n gives

k = −n tan(α/2)

(Figure 4). k is a pure imaginary quaternion, that is, a
three-dimensional vector.

\begin{itemize}
\tightlist
\item
  Magnitude: \textbar k\textbar{} = \textbar q₁ − q₂\textbar{} /
  \textbar q₁ + q₂\textbar{} = tan(α/2). Since the norm is
  multiplicative, this is the same formula as for complex numbers.
\item
  Direction: n is the axis representing in which direction the two
  states are displaced.
\end{itemize}

The distance map is x = k⁻¹ = n cot(α/2). From the relation between two
states, a three-dimensional point with direction and distance is read
out.

\begin{figure}
\centering
\pandocbounded{\includegraphics[keepaspectratio,alt={Extension to quaternions. Projecting the relative phase u on S³ from the antipode of the observer gives the three-dimensional vector k (left; S³ is shown as S² with one dimension suppressed). In the left and right frames the length is common and the direction differs (right).}]{figs/fig04_quaternion_S3_to_R3_en.png}}
\caption{Extension to quaternions. Projecting the relative phase u on S³
from the antipode of the observer gives the three-dimensional vector k
(left; S³ is shown as S² with one dimension suppressed). In the left and
right frames the length is common and the direction differs (right).}
\end{figure}

\subsubsection{4.2 Relation to the central projection of the previous
paper}\label{relation-to-the-central-projection-of-the-previous-paper}

The previous paper {[}S6{]} showed that the polar form of a quaternion q
= \textbar q\textbar(cos θ + n sin θ) is exactly the central projection
onto the tangent hyperplane at the north pole, u\_g = v/q₀ = n tan θ.
Applying this to the square root of the relative phase p = e\^{}(nα/2)
(u = p²), its image is n tan(α/2) = −k.

\textbf{Proposition 1.} The coupling map k is the stereographic
projection of the relative phase u and, at the same time, the central
projection of the square root of the relative phase p = e\^{}(nα/2).

The coupling strength becomes \textbar k\textbar{} = 1 when α = π/2
(orthogonal), and then p = e\^{}(nπ/4) lies on the light cone in the
sense of the previous paper (θ = 45°, the real part of p² is 0).

\textbf{Corollary.} A relation with coupling coefficient 1 is a relation
whose square root of the relative phase lies on the light cone. In the
most symmetric configuration of the previous paper, q = ±e\^{}(πn/4),
each element sits exactly at this position.

\subsubsection{4.3 Left and right frames}\label{left-and-right-frames}

If the relative phase is defined by right division v = q₂q₁⁻¹, then v =
q₁ u q₁⁻¹, and the coupling coefficient becomes k\_R = q₁ k\_L q₁⁻¹. The
magnitude \textbar k\textbar{} is common, and the direction is only
rotated by conjugation with q₁. The strength of the coupling does not
depend on the observing frame; the direction depends on the frame. This
is the same structure as, in the visual angle of the triangle in the
previous paper, n\_L and n\_R being rotated separately by left and right
multiplication.

\subsubsection{4.4 When the amplitudes differ: the 1+3
decomposition}\label{when-the-amplitudes-differ-the-13-decomposition}

Writing u = ρ e\^{}(nα), and using that (1 + u)⁻¹ = (1 + ū)/\textbar1 +
u\textbar² commutes with (1 − u),

k = {[}(1 − ρ²) − 2 Im u{]} / \textbar1 + u\textbar².

The scalar part is the amplitude mismatch, and the vector part is the
phase coupling with an orientation. The complex ``real part = mismatch,
imaginary part = phase'' extends directly to a 1+3 decomposition.

\subsubsection{4.5 Invariance and composition
law}\label{invariance-and-composition-law}

\textbar k\textbar{} is invariant under q ↦ a q b (four-dimensional
rotations, a, b unit quaternions) and a common real factor. k itself is
invariant under common multiplication from the left, and rotates as a
vector under common multiplication from the right. If the axes are
aligned, the composition law is the same as in §3.4, but if the axes are
misaligned, the composition depends on the order of multiplication and
cross terms of cross-product type appear.

\begin{center}\rule{0.5\linewidth}{0.5pt}\end{center}

\subsection{5. The localized observer and local
flatness}\label{the-localized-observer-and-local-flatness}

\subsubsection{5.1 What cannot be
distinguished}\label{what-cannot-be-distinguished}

A localized observer q\_o reads only the relative phases u\_j =
q\_o⁻¹q\_j with itself as reference. Itself is u = 1, so in the coupling
map it comes to the origin. The coupling map is the stereographic
projection centred on the observer, and stereographic projection
preserves angles. Therefore an observer who measures only by visual
angles cannot distinguish the following.

\begin{enumerate}
\def\labelenumi{\arabic{enumi}.}
\tightlist
\item
  The uniform curvature of S³ and flat ℝ³ (S³ is also locally flat)
\item
  The coupling map k and the distance map x = 1/k (inversion preserves
  angles)
\item
  The choice of left or right frame (the observer has only its own
  frame)
\item
  The global scale (distance is obtained only as a ratio to a size)
\end{enumerate}

What the observer obtains are the directions n\_j to other states, the
angles between directions, and ratios of distances; this is complete as
the information of a three-dimensional space identified up to
angle-preserving maps.

\subsubsection{5.2 What is visible: tides}\label{what-is-visible-tides}

What cannot be distinguished is only the uniform part of curvature.
Direction-dependent distortion, that is, the component by which the
shape of the set of states in the observer's field of view is squeezed
or stretched in one direction, is visible as a deformation of ratios
between angles. This is the same structure as the equivalence principle:
a uniform gravitational field can be removed locally, but tides cannot.

\subsubsection{5.3 Local flatness is a property of the
readout}\label{local-flatness-is-a-property-of-the-readout}

To the question of §1.1, this paper answers as follows.

\textbf{Local flatness is not an approximation but a property of the
observer-centred conformal readout. Curvature appears only as defects
(deficit angles at the seams, §9).}

In the equal-amplitude case, the basis of the readout is the round S³
projected stereographically, so however the states are distributed, the
geometry of the field of view is conformally flat and tides do not
appear in principle. Equal amplitude corresponds to a uniform state
without gravity. Tides appear when the amplitudes are not uniform, and
the scalar part of §4.4 enters the readout as a factor that stretches
and shrinks distances direction by direction (§6.3).

\begin{center}\rule{0.5\linewidth}{0.5pt}\end{center}

\subsection{6. Standing waves and mass
(hypothesis)}\label{standing-waves-and-mass-hypothesis}

\subsubsection{6.1 No travelling phase}\label{no-travelling-phase}

We consider that quaternions have no travelling phase and that states
are standing waves. Under this premise, the relative phase α of two
states does not flow in time, and the readout of distance is statically
fixed.

\subsubsection{6.2 Mass is the norm}\label{mass-is-the-norm}

For a single quaternion, the only quantity invariant under the
four-dimensional rotations q ↦ a q b is the norm \textbar q\textbar².
This is the Euclidean version of the structure in which the
Lorentz-invariant norm of the four-momentum gives the rest mass. What
reaches the observer is only the amplitude ratio ρ =
\textbar q\_j\textbar{} / \textbar q\_o\textbar, and the global scale
cannot be read.

A standing wave is a state in which counter-propagating components
balance and the net momentum is zero, that is, the rest frame itself.
Just as a box enclosing light has inertial mass as a whole, the energy
of a standing wave is proportional to the square of the amplitude. We
therefore place the following hypothesis.

\textbf{Hypothesis 1.} Mass is proportional to the squared norm
\textbar q\textbar², and is read out statically from the observer as the
amplitude ratio ρ.

Since gravitational mass and inertial mass come from the same quantity,
the equivalence principle is built into the premise. Reading out a mass
ratio requires no clock; it requires only the state that serves as the
reference for the ratio. Mass is the value obtained by reading out the
internal state from outside, and the conversion ratio between the
internal and external scales is not derived in this paper.

\subsubsection{6.3 The source of gravity is the non-uniformity of the
norm}\label{the-source-of-gravity-is-the-non-uniformity-of-the-norm}

Since the global scale is not visible, what produces gravity is not
\textbar q\textbar{} itself but its non-uniformity. Regarding Φ ∝ ln
\textbar q\textbar{} as a potential, the mass density corresponds to its
Laplacian, and the mass in an enclosed region to the flux of its
gradient. In gravitational lensing, the isotropic magnification
(convergence) is proportional to the projected mass density, which is
consistent with the scalar part of k producing isotropic distortion.

\subsubsection{6.4 Motion is a left--right
imbalance}\label{motion-is-a-leftright-imbalance}

Four-dimensional rotations divide into rotations from the left and
rotations from the right. The previous paper {[}S6{]} showed that what
preserves the true-square closure Σq² = 0 is only the rotation q ↦ u q
u⁻¹ multiplying the same unit from both sides, and real scaling, and
that one-sided multiplication breaks the closure.

\textbf{Hypothesis 2.} Rest is the state in which left and right are
balanced (rotating with the same unit from both sides, preserving the
closure), and motion is the breaking of the left--right balance
(one-sided multiplication, breaking the closure). The travelling phase
is not ``absent''; it is balanced to zero in the rest state and appears
with motion.

Under this hypothesis, instead of the equivalence principle holding
automatically, it is replaced by the dynamical question ``does the
gradient of the norm (gravity) cause a left--right imbalance (motion)?''
If this holds, free fall can be read as a phenomenon in which the
non-uniformity of amplitude tilts a standing wave toward the
travelling-wave side.

\begin{center}\rule{0.5\linewidth}{0.5pt}\end{center}

\subsection{7. Anonymity: reading out connectivity from an anonymous
point
set}\label{anonymity-reading-out-connectivity-from-an-anonymous-point-set}

\subsubsection{7.1 The problem with simplicial
methods}\label{the-problem-with-simplicial-methods}

In a simplicial method that builds structure from a complete network of
vertices and edges, vertices have identities, and the connectivity
information of which vertex is joined to which is supplied from outside.
This is contrary to anonymity. What may be given is only a collection of
anonymous states, and connectivity must be read out from it.

\subsubsection{7.2 Readout by the convex
hull}\label{readout-by-the-convex-hull}

Equal-amplitude quaternions are points on S³ ⊂ ℝ⁴. Taking the convex
hull of a collection of anonymous points, its faces are tetrahedra, and
a truss filling S³ with tetrahedra is obtained. The convex hull does not
depend at all on the ordering or numbering of the points.

When projected stereographically from the antipode of the observer q\_o,
the faces of the convex hull whose caps do not contain the pole of the
projection coincide with the Delaunay triangulation of the projected
points (the correspondence between the convex hull of a point set on a
sphere and the Delaunay triangulation {[}X6{]}). Figure 5 confirms
numerically, for 30 anonymous points on S² as a one-dimension-lower
analogue of S³, that the projections of the faces of the convex hull
coincide with the Delaunay triangulation of the projected points (53
triangles).

\begin{figure}
\centering
\pandocbounded{\includegraphics[keepaspectratio,alt={The convex hull of an anonymous point set (left) and the Delaunay triangulation by observer-centred stereographic projection (right). Red is an example with no other point inside the circumcircle.}]{figs/fig05_convex_hull_delaunay_en.png}}
\caption{The convex hull of an anonymous point set (left) and the
Delaunay triangulation by observer-centred stereographic projection
(right). Red is an example with no other point inside the circumcircle.}
\end{figure}

\subsubsection{7.3 Consistency with the readouts so
far}\label{consistency-with-the-readouts-so-far}

\begin{enumerate}
\def\labelenumi{\arabic{enumi}.}
\tightlist
\item
  \textbf{Independent of the choice of map.} The Delaunay test (no other
  point inside the circumsphere of a tetrahedron) is invariant under
  maps that take spheres to spheres. Whichever of the maps k and x,
  indistinguishable for a localized observer, is used for the test, the
  same truss is obtained.
\item
  \textbf{Amplitude enters as a weight.} When amplitudes differ, the
  points are displaced radially from S³, and their convex hull
  generalizes to a weighted Delaunay triangulation (regular
  triangulation) {[}X7{]}. Points with larger amplitude (mass,
  Hypothesis 1) change the surrounding connectivity more. The specific
  correspondence between amplitude and weight is not derived in this
  paper.
\item
  \textbf{Can be assembled by local rules.} The circumsphere test can be
  performed with nearby points only, and places that do not satisfy the
  condition can be repaired by local rearrangements (flips) {[}X7{]}.
  Without an eye that surveys the whole, the truss is assembled and
  rearranged by local operations alone.
\end{enumerate}

\subsubsection{7.4 Degeneracy arising from
anonymity}\label{degeneracy-arising-from-anonymity}

When five or more points lie on the same sphere, the triangulation is
not determined uniquely. Also, there are point sets with different
structures even though their sets of unlabelled distances are the same.
Giving names from outside would resolve this degeneracy, but that would
break anonymity. This paper does not regard this degeneracy as a defect
to be eliminated, and places it as the principle of the next section.

\begin{center}\rule{0.5\linewidth}{0.5pt}\end{center}

\subsection{8. Partial projection and
indeterminacy}\label{partial-projection-and-indeterminacy}

\subsubsection{8.1 The rigidity count}\label{the-rigidity-count}

For a three-dimensional truss joining V points with E members, Maxwell's
count and Calladine's relation {[}X8, X9{]}

F − S = 3V − 6 − E

hold. F is the number of degrees of freedom of deformation excluding
rigid-body motion (soft modes), and S is the number of internal stresses
the members hold by constraining one another (self-stresses).

\subsubsection{8.2 Three phases of
projection}\label{three-phases-of-projection}

{\def\LTcaptype{none} % do not increment counter
\begin{longtable}[]{@{}
  >{\raggedright\arraybackslash}p{(\linewidth - 4\tabcolsep) * \real{0.3333}}
  >{\raggedright\arraybackslash}p{(\linewidth - 4\tabcolsep) * \real{0.3333}}
  >{\raggedright\arraybackslash}p{(\linewidth - 4\tabcolsep) * \real{0.3333}}@{}}
\toprule\noalign{}
\begin{minipage}[b]{\linewidth}\raggedright
Phase
\end{minipage} & \begin{minipage}[b]{\linewidth}\raggedright
Members
\end{minipage} & \begin{minipage}[b]{\linewidth}\raggedright
Result
\end{minipage} \\
\midrule\noalign{}
\endhead
\bottomrule\noalign{}
\endlastfoot
Zero projection & E = 0 & F is maximal. No constraint; no structure
forms \\
Partial projection & E ≈ 3V − 6 & Both F and S remain moderately.
Skeleton and room for rearrangement coexist \\
Full projection & E = V(V−1)/2 & F = 0, S is huge. It becomes a rigid
crystal and does not evolve \\
\end{longtable}
}

Figure 6 shows the three phases in a two-dimensional analogue (2V − 3).
The example of partial projection is an isostatic truss grown from a
seed by Henneberg-type operations (joining a new point to two existing
points) {[}X10{]}.

\begin{figure}
\centering
\pandocbounded{\includegraphics[keepaspectratio,alt={Three phases of projection. From the left: zero projection, partial projection, full projection.}]{figs/fig06_three_phases_rigidity_en.png}}
\caption{Three phases of projection. From the left: zero projection,
partial projection, full projection.}
\end{figure}

The previous paper and the catalogue {[}S3, S6{]} stated that, as long
as spacetime is read stably, the map must be at least a partial
projection: with zero projection spacetime does not appear, and with
full projection no unreadable interior, indeterminacy or discreteness
remains. This section makes this concrete in the language of trusses.

\subsubsection{8.3 Indeterminacy and time}\label{indeterminacy-and-time}

\begin{itemize}
\tightlist
\item
  \textbf{Indeterminacy:} The soft modes F and the non-uniqueness of the
  Delaunay triangulation (§7.4) carry the indeterminacy of the
  structure.
\item
  \textbf{Time:} It appears as a sequence of local rearrangements
  (flips). Since standing waves have no travelling phase (§6.1), time
  arises not from the advance of phase but from the rearrangement of
  structure.
\item
  \textbf{Complementarity:} F − S is fixed by the numbers of points and
  members alone. Increasing freedom reduces constraint, increasing
  constraint reduces freedom, and both cannot be fixed at the same time.
\end{itemize}

\subsubsection{8.4 Composition requires a
process}\label{composition-requires-a-process}

Since the structure is not determined uniquely, composition is not
``solve the equations and obtain one answer''; it can only be defined as
a process with a procedure and a history. The result depends on where,
and in what order, the degeneracy was resolved. Among the many possible
structures, the actual history selects one. This history dependence can
be the origin of the arrow of time.

In glass science, it is known that network glasses form most easily near
the isostatic point where the number of constraints balances the degrees
of freedom {[}X11, X12{]}. The proportion of partial projection must
also be decided not from outside but by the structure itself. The
mechanism is described in the next section.

\begin{center}\rule{0.5\linewidth}{0.5pt}\end{center}

\subsection{9. Supercooling, frustration and
defects}\label{supercooling-frustration-and-defects}

\subsubsection{9.1 Supercooling and seeds}\label{supercooling-and-seeds}

\begin{itemize}
\tightlist
\item
  \textbf{Supercooled state:} A highly symmetric, degenerate metastable
  state in which all relations exist potentially but none has yet been
  fixed as a member.
\item
  \textbf{Seed:} The first local block (one tetrahedron) in which the
  degeneracy is resolved. To preserve anonymity, it is not placed from
  outside but arises spontaneously from fluctuations of the
  configuration.
\item
  \textbf{Crystallization:} A fixed tetrahedron resolves the degeneracy
  of neighbouring faces, and the front of fixing spreads (§10).
\end{itemize}

This is the same structure as supercooling by icosahedral local order
{[}X13{]}.

\subsubsection{9.2 Tetrahedral
frustration}\label{tetrahedral-frustration}

The dihedral angle of the regular tetrahedron is arccos(1/3) ≈ 70.53°;
placing five around one edge gives 5 × 70.53° = 352.64°, leaving a gap
of 7.36°. Regular tetrahedra cannot fill flat three-dimensional space
without gaps (Figure 7, left).

In S³, however, there exists a regular polytope that 600 regular
tetrahedra fill completely (the 600-cell). Five cells surround each
edge, and the dihedral angle of the cells is exactly 72°. Its 120
vertices are unit quaternions and form the binary icosahedral group
{[}X3, X5{]} (Figure 7, right).

\begin{figure}
\centering
\pandocbounded{\includegraphics[keepaspectratio,alt={View along a shared edge. In flat space a gap remains (left); in S³ it closes (right).}]{figs/fig07_tetrahedral_frustration_600cell_en.png}}
\caption{View along a shared edge. In flat space a gap remains (left);
in S³ it closes (right).}
\end{figure}

\subsubsection{9.3 Partial projection is enforced
geometrically}\label{partial-projection-is-enforced-geometrically}

This is exactly the structure known in the physics of supercooled
liquids {[}X14, X15{]}. Perfect icosahedral order is realized only in
the curved space S³, and mapping it to flat space necessarily introduces
defect lines (disclinations). In the context of this paper:

\begin{enumerate}
\def\labelenumi{\arabic{enumi}.}
\tightlist
\item
  On S³ (the space of unit quaternions), crystallization tends toward
  perfect order.
\item
  Since the map of a localized observer is locally flat (§5), that order
  cannot be mapped completely, and defect lines necessarily remain.
\item
  Defect lines have deficit angles and become the carriers of curvature
  (tides).
\item
  When the number of states greatly exceeds 120, complete filling is
  impossible even on S³, so frustration is essentially unavoidable.
\end{enumerate}

\textbf{Proposition 2 (hypothesis).} The proportion of partial
projection is not given by an external parameter but is enforced
geometrically by tetrahedral frustration. The system cannot become a
fully projected rigid crystal, and room for defects and rearrangement
always remains.

In the methods of condensed-matter physics, curved space is introduced
as a tool to describe order. In this paper, the state space itself is
S³, and curved space is the starting point.

\begin{center}\rule{0.5\linewidth}{0.5pt}\end{center}

\subsection{10. Geometric growth and
inflation}\label{geometric-growth-and-inflation}

\subsubsection{10.1 Counting by the number of
processes}\label{counting-by-the-number-of-processes}

Since space is not presupposed, growth is counted not as a front
advancing at constant speed in an existing space but by the number of
processes. In one process, tetrahedra attach simultaneously to all
exposed faces. When one tetrahedron attaches, one face is closed and
three new faces are exposed. Letting r (r ≤ 3) be the factor per process
after subtracting the part closed by new tetrahedra sharing faces with
one another, after n processes the number of bondable faces and the
number of blocks are both proportional to rⁿ (Figure 8, left).

Geometric growth is the default behaviour of this process; what is
special is not the growth but the decrease of the factor as faces close.

\begin{figure}
\centering
\pandocbounded{\includegraphics[keepaspectratio,alt={The tree of faces (left), and the growth of the number of blocks for each factor r (right). The dashed line is an illustrative model including closure and is not derived.}]{figs/fig08_geometric_growth_efold_en.png}}
\caption{The tree of faces (left), and the growth of the number of
blocks for each factor r (right). The dashed line is an illustrative
model including closure and is not derived.}
\end{figure}

\subsubsection{10.2 Number of e-folds}\label{number-of-e-folds}

Taking one block as the unit volume, V ∝ rⁿ, and the length scale is a ∝
V\^{}(1/3), so the number of e-folds is

N\_e = (n/3) ln r.

For r = 3, the number of processes needed for 60 e-folds is n = 180/ln 3
≈ 164, and then V ≈ 10⁷⁸. In standard inflation theory the expansion
rate is tuned by the inflaton potential; here the factor is fixed by
geometry, namely the number of faces of a tetrahedron.

\subsubsection{10.3 How it ends, and how it looks to the
observer}\label{how-it-ends-and-how-it-looks-to-the-observer}

\begin{itemize}
\tightlist
\item
  \textbf{End:} When faces begin to close through the resolution of
  degeneracy or the curvature of S³, r approaches 1 and growth stops.
  The end of inflation is given structurally as a change of topology.
\item
  \textbf{Appearance:} In the growth stage, the number of faces and the
  number of blocks both increase as rⁿ, so the ratio of surface area to
  volume is kept constant. This is a characteristic not of flat space
  (where the ratio is inversely proportional to the radius) but of
  hyperbolic geometry. Since the global scale is not visible to a
  localized observer (§5.1), the increase of the number of blocks around
  it by a constant factor should be read out as exponential expansion
  (hypothesis).
\end{itemize}

\subsubsection{10.4 Correspondence with standard
theory}\label{correspondence-with-standard-theory}

{\def\LTcaptype{none} % do not increment counter
\begin{longtable}[]{@{}
  >{\raggedright\arraybackslash}p{(\linewidth - 4\tabcolsep) * \real{0.3333}}
  >{\raggedright\arraybackslash}p{(\linewidth - 4\tabcolsep) * \real{0.3333}}
  >{\raggedright\arraybackslash}p{(\linewidth - 4\tabcolsep) * \real{0.3333}}@{}}
\toprule\noalign{}
\begin{minipage}[b]{\linewidth}\raggedright
\end{minipage} & \begin{minipage}[b]{\linewidth}\raggedright
Standard inflation
\end{minipage} & \begin{minipage}[b]{\linewidth}\raggedright
This paper
\end{minipage} \\
\midrule\noalign{}
\endhead
\bottomrule\noalign{}
\endlastfoot
What grows exponentially & The scale of existing space & The number of
fixed relations (blocks) \\
Driver & The potential of the inflaton field & Resolution of degeneracy
(default behaviour) \\
End & The field rolls down the potential & Closure of faces, filling of
S³, frustration \\
Reheating & Decay of the inflaton & Release of the latent heat of
crystallization \\
Uniformity & Stretching of a causally connected small region & A chain
from a single seed \\
Flatness & Flattening by stretching & Uniform curvature is invisible to
a localized observer \\
Primordial fluctuations & Quantum fluctuations of the inflaton & Frozen
defect lines and non-uniqueness of rearrangement \\
\end{longtable}
}

Guth's original inflationary theory {[}X17{]} was a picture of a
supercooled false vacuum and nucleation of bubbles, but it had the
problem that the expansion of space was faster than the growth of
bubbles, so bubbles did not connect and inflation could not end
{[}X18{]}, and it was replaced by the slow-roll type {[}X19{]}. In the
picture of this paper, space itself is not presupposed and what expands
is the amount of fixed structure, so this problem is avoided
structurally in the form ``it stops when faces close.''

Topological defects produced at phase transitions are connected, as the
Kibble mechanism {[}X20{]} and its condensed-matter version the
Kibble--Zurek mechanism {[}X21{]}, by a scaling law in which the defect
density is determined by the speed of the phase transition, and this has
been verified experimentally {[}X22{]}. It provides a foothold for
checking the correspondence between the defect lines of this paper and
primordial fluctuations.

\subsubsection{10.5 Relation to task 3 of the previous
paper}\label{relation-to-task-3-of-the-previous-paper}

Task 3 of the previous paper {[}S6{]} was the dynamical question of
whether the all-to-all self-interaction of highly symmetric quaternions
produces inflation-like expansion. This paper answers it from the side
of structure (the counting of faces), as a stage before writing down the
self-interaction. Whether the two are different expressions of the same
phenomenon remains a task to be checked after implementing the
self-interaction.

\begin{center}\rule{0.5\linewidth}{0.5pt}\end{center}

\subsection{11. A bridge to octonions: the spacetime equation is the
null condition of split
octonions}\label{a-bridge-to-octonions-the-spacetime-equation-is-the-null-condition-of-split-octonions}

\subsubsection{11.1 Correspondence}\label{correspondence}

Take the spacetime-side quaternion a = t + x i + y j + z k and the
internal-side quaternion b = R + Q₁ i + Q₂ j + Q₃ k, and consider the
Cayley--Dickson construction o = (a, b) with product (a, b)(c, d) = (ac
+ γ d̄ b, da + b c̄) and conjugate (a, b)* = (ā, −b). The norm is N(o) =
\textbar a\textbar² − γ\textbar b\textbar²; γ = −1 gives the ordinary
octonions and γ = +1 the split octonions (signature (4,4)) {[}X4{]}. For
the split octonions

x² + y² + z² + t² − (R² + Q₁² + Q₂² + Q₃²) = 0　⇔　N(o) =
0　⇔　\textbar a\textbar{} = \textbar b\textbar,

and if a ≠ 0 the relative phase u = a⁻¹b is a unit quaternion (a point
on S³) (Figure 9).

\begin{figure}
\centering
\pandocbounded{\includegraphics[keepaspectratio,alt={The spacetime equation is the null condition of split octonions and is equivalent to a pair of equal-amplitude quaternions.}]{figs/fig09_split_octonion_null_en.png}}
\caption{The spacetime equation is the null condition of split octonions
and is equivalent to a pair of equal-amplitude quaternions.}
\end{figure}

\subsubsection{11.2 What can be read}\label{what-can-be-read}

\begin{enumerate}
\def\labelenumi{\arabic{enumi}.}
\tightlist
\item
  \textbf{The spacetime equation is a pair of equal-amplitude
  quaternions.} The structure of this paper, ``reading out relations
  from the relative phase of two equal-amplitude quaternions,'' is the
  spacetime equation itself.
\item
  \textbf{The deviation from null is the amplitude mismatch.} When
  \textbar a\textbar{} ≠ \textbar b\textbar, a scalar part appears in
  the coupling coefficient (§4.4) and becomes a candidate for
  gravitational distortion (§6.3).
\item
  \textbf{Zero closure requires the split form.} The norm of the
  ordinary octonions is positive definite and has no nontrivial null
  elements. Requiring nontrivial solutions selects the split form. This
  is the same logic by which Σxₙ² = 0 has only the trivial solution over
  the reals and requires complex numbers (Stage ①).
\item
  \textbf{The internal side appears only through the relative phase and
  the norm ratio.} What is visible from the spacetime side is only u and
  the norm ratio; the state of the internal side itself is not read out.
\end{enumerate}

The previous paper {[}S6{]} showed that the true square of an octonion
repeats the same form o² = (q² − \textbar p\textbar²) + 2q₀ p \(\ell\).
The norm of the split form in this section is of the form on the side of
the norm using conjugation, and its relation to the side of the true
square remains a task. How non-associativity appears in chains of
relations (the closure of triangles) is also unresolved.

\begin{center}\rule{0.5\linewidth}{0.5pt}\end{center}

\subsection{12. Comparison with previous
research}\label{comparison-with-previous-research}

\subsubsection{12.1 The closest previous
research}\label{the-closest-previous-research}

\begin{itemize}
\tightlist
\item
  \textbf{Quantum Graphity} {[}X1{]} placed degrees of freedom on a
  complete graph and considered, from a model invariant under
  permutations of points, the process by which geometry arises through a
  phase transition by cooling (geometrogenesis). Starting from a
  complete graph and imposing permutation invariance is shared with this
  paper. On the other hand, the vertices carry labels, and the phase
  transition requires an external heat bath. This paper gives no names
  to vertices, places neither a heat bath nor a Hamiltonian, and reads
  out connectivity from the convex hull of an anonymous point set.
\item
  \textbf{The Selection--Stitch Model} {[}X2{]} attempts to explain
  inflation and reheating as a process in which the tensor network of a
  discrete vacuum crystallizes through a frustrated tetrahedral foam.
  The combination of tetrahedral frustration and inflation by
  crystallization is shared with this paper. On the other hand, it
  presupposes a network and a coordination number, and selects phases by
  energy minimization. This paper starts from the fact that the state
  space itself is S³ (unit quaternions), does not supply connectivity
  from outside, and counts the growth factor by the number of faces.
\end{itemize}

\subsubsection{12.2 Condensed-matter physics and glass
science}\label{condensed-matter-physics-and-glass-science}

The method of placing a template of perfect order in curved space and
introducing defects by flattening {[}X14, X15{]}, constraint theory and
rigidity percolation {[}X11, X12{]}, and the kinetics of nucleation and
growth {[}X16{]} are direct precedents for the structure of §8--§10 of
this paper. The novelty of this paper lies in placing these in the
problem of reading out spacetime and treating curved space not as a tool
but as the state space.

\subsubsection{12.3 Discrete spacetimes that supply connectivity from
outside}\label{discrete-spacetimes-that-supply-connectivity-from-outside}

Curvature as a deficit angle {[}X23{]} and causal dynamical
triangulations, in which four dimensions emerge from gluing simplices
{[}X24{]}, share with this paper the point that curvature concentrates
at seams, but they supply the rules of connectivity from outside.

\subsubsection{12.4 Novelty of this paper}\label{novelty-of-this-paper}

\begin{enumerate}
\def\labelenumi{\arabic{enumi}.}
\tightlist
\item
  Coupling and distance are read out from the relative phase of
  quaternions as an observer-centred stereographic projection (§3--§4).
  The coupling map is the central projection of the square root of the
  relative phase, and coupling coefficient 1 corresponds to the light
  cone (Proposition 1).
\item
  Connectivity is read out from the convex hull of an anonymous point
  set (§7).
\item
  Non-uniqueness is placed as a principle, partial projection =
  indeterminacy, and composition is defined as a process (§8).
\item
  Without a heat bath or a Hamiltonian, geometric growth and the number
  of e-folds are fixed by the factor of faces, and the end is given
  structurally (§10).
\item
  The state space itself is S³, and curved space is the starting point
  (§9).
\item
  The spacetime equation is the null condition of split octonions and is
  equivalent to a pair of equal-amplitude quaternions (§11).
\end{enumerate}

\begin{center}\rule{0.5\linewidth}{0.5pt}\end{center}

\subsection{13. Verification plan and falsification
conditions}\label{verification-plan-and-falsification-conditions}

The claims of this paper divide into mathematically identical ones,
structural analogies, and hypotheses. The classification is as follows.

{\def\LTcaptype{none} % do not increment counter
\begin{longtable}[]{@{}
  >{\raggedright\arraybackslash}p{(\linewidth - 2\tabcolsep) * \real{0.5000}}
  >{\raggedright\arraybackslash}p{(\linewidth - 2\tabcolsep) * \real{0.5000}}@{}}
\toprule\noalign{}
\begin{minipage}[b]{\linewidth}\raggedright
Claim
\end{minipage} & \begin{minipage}[b]{\linewidth}\raggedright
Category
\end{minipage} \\
\midrule\noalign{}
\endhead
\bottomrule\noalign{}
\endlastfoot
Angular-size distance and coupling coefficient (§3), the quaternion
coupling coefficient and Proposition 1 (§4), the correspondence between
convex hull and Delaunay triangulation (§7.2), the rigidity count
(§8.1), the tetrahedral gap and the 600-cell (§9.2), the null condition
of split octonions (§11.1) & Mathematically identical \\
Correspondence with the reflection coefficient (§3.2), correspondence
between supercooling and inflation (§10.4) & Structural analogy \\
Hypotheses 1 and 2 (§6), the source of tides (§5.2, §6.3), Proposition 2
(§9.3), readout as exponential expansion (§10.3) & Hypothesis \\
\end{longtable}
}

Numerical experiments to test the hypotheses are planned as follows.

\begin{itemize}
\tightlist
\item
  \textbf{V1 Four-dimensional version of the convex-hull readout:} Take
  the convex hull in ℝ⁴ of N anonymous unit quaternions, and confirm
  that the observer-centred stereographic projection coincides with the
  three-dimensional Delaunay triangulation. Record the distributions of
  the number of edges per vertex and the number of tetrahedra per edge.
\item
  \textbf{V2 Proportion of partial projection:} Project members with a
  degeneracy criterion (making members only where the triangulation is
  determined uniquely), and measure F, S and E/(3V − 6) while varying N.
\item
  \textbf{V3 Supercooling and seeds:} Take as the supercooled state a
  highly symmetric configuration (the vertices of the 600-cell, the most
  symmetric configuration of the previous paper, etc.) with a small
  disorder added, and in the process of attaching tetrahedra by local
  circumsphere tests, measure the factor r\_n per process and N\_e.
\item
  \textbf{V4 Defects:} Measure the number of tetrahedra surrounding an
  edge (ideally 5) and the distribution of deficit angles in the
  observer's map.
\item
  \textbf{V5 Source of tides:} Compare the deviation (shear) of the
  observer's field-of-view map from conformal flatness between
  equal-amplitude configurations and configurations with non-uniform
  amplitude.
\end{itemize}

The falsification conditions are as follows.

\begin{enumerate}
\def\labelenumi{\arabic{enumi}.}
\tightlist
\item
  If no stage in which the factor r\_n clearly exceeds 1 appears in V3,
  the claim of geometric growth in §10 is rejected.
\item
  If E/(3V − 6) does not settle to a value independent of N in V2,
  Proposition 2 (the proportion of partial projection is determined
  geometrically) is rejected.
\item
  If shear appears in equal-amplitude configurations in V5, ``with equal
  amplitude, tides do not appear in principle'' in §5.3 is wrong.
\item
  If defect lines do not remain after flattening in V4, §9.3 is
  rejected.
\end{enumerate}

\begin{center}\rule{0.5\linewidth}{0.5pt}\end{center}

\subsection{14. Conclusion}\label{conclusion}

A readout that takes the visual angle as distance was defined from the
relation between two states. For equal-amplitude complex numbers we
obtained the coupling coefficient g = tan(Δθ/2) and the angular-size
distance D = cot(Δθ/2); in phase is at infinity and antiphase is
contact. For quaternions the coupling coefficient becomes the
three-dimensional vector k = −n tan(α/2), and a three-dimensional point
with direction and distance is read out from the relation between two
states alone. k is the stereographic projection of the relative phase
and also the central projection of the square root of the relative
phase, and coupling coefficient 1 corresponds to the light cone.

Since this readout preserves angles, a localized observer cannot
distinguish uniform curvature, scale, map, or left/right frame, and
local flatness is a property of the readout. Only tides are visible, and
we consider their source to be the non-uniformity of amplitude.

Connectivity is read out from the convex hull of an anonymous point set.
Placing the non-uniqueness of the triangulation as a principle, only
partial projection, between zero projection and full projection,
combines a skeleton with room for rearrangement, and composition
requires a process. When a seed enters a supercooled, highly symmetric
state, geometric crystallization in which the number of faces multiplies
by r per process occurs, giving the number of e-folds N\_e = (n/3) ln r,
and it ends when faces close. Regular tetrahedra cannot fill flat space
but close in S³, so defect lines necessarily remain in a locally flat
readout, and partial projection is enforced without external parameters.

Finally, it was shown that the author's spacetime equation is the null
condition of split octonions and is equivalent to the quaternions of the
spacetime side and the internal side having equal amplitude. The
relational readout of quaternions in this paper is a bridge to the stage
of octonions.

\begin{center}\rule{0.5\linewidth}{0.5pt}\end{center}

\subsection{References}\label{references}

\subsubsection{Self-citations (Concept
DOI)}\label{self-citations-concept-doi}

{[}S1{]} N. Kihara, ``Geometric Formulation of 4-Dimensional Space via
Central Projection,'' 2026-04-05. Concept DOI 10.5281/zenodo.19427779.
{[}S2{]} N. Kihara, ``Anonymous Equal-Amplitude Composite-Wave Model:
Basic Axiom System v9,'' 2026-07-24. Concept DOI
10.5281/zenodo.21315735. {[}S3{]} N. Kihara, ``A Catalogue of Symmetries
and Magic Numbers for Reading Out Spacetime from a State Space: The
Answer-Check Standard of a Research Programme Taking Anonymity and
Self-Consistency as Axioms,'' 2026-10-18. Concept DOI
10.5281/zenodo.23444451. {[}S4{]} N. Kihara, ``Reading Out Time, Space,
Mass, Charge, Translational Velocity and Spin from the Relations and
Interference of Complex Waves without a Background Spacetime: An
Integrated Hypothesis in the Order of the Thought Experiment, from
Residents of a Closed Circle to Residents of a Four-Dimensional
Spherical Shell, and on to α, the Four-Dimensional Lattice and E8,''
2026-10-15. Concept DOI 10.5281/zenodo.23384928. {[}S5{]} N. Kihara,
``Dynamics of a Self-Consistent Relational-Wave Closed System: From the
Local Instability of the 90-Degree Relative Equilibrium to Inflation,
Spontaneous Symmetry Breaking, Grassmann Geometry, Terminal Universality
and Relationalism,'' 2026-09-12. Concept DOI 10.5281/zenodo.22729146.
{[}S6{]} N. Kihara, ``A Map from a Set of Complex Numbers to Spacetime
(1+3): Readout by Quaternions and Partial Projection,'' 2026-10-18.
Concept DOI 10.5281/zenodo.23451494.

\subsubsection{External references}\label{external-references}

{[}X1{]} T. Konopka, F. Markopoulou, S. Severini, ``Quantum graphity: a
model of emergent locality,'' Phys. Rev.~D 77, 104029 (2008); the
original is T. Konopka, F. Markopoulou, L. Smolin, ``Quantum Graphity,''
arXiv:hep-th/0611197 (2006). {[}X2{]} R. Kulkarni, ``Emergence of cosmic
flatness, inflation, and reheating from topological phase transitions in
a discrete vacuum,'' Zenodo (2026). Concept DOI 10.5281/zenodo.18727238
(the first version of 2026-02-21 was titled ``Geometric Phase
Transitions in a Discrete Vacuum: Deriving Cosmic Flatness, Inflation,
and Reheating from Tensor Network Topology''). {[}X3{]} J. H. Conway, D.
A. Smith, \emph{On Quaternions and Octonions: Their Geometry,
Arithmetic, and Symmetry} (A K Peters, Natick, 2003). {[}X4{]} J. C.
Baez, ``The Octonions,'' Bull. Amer. Math. Soc. 39, 145--205 (2002).
{[}X5{]} H. S. M. Coxeter, \emph{Regular Polytopes}, 3rd ed.~(Dover, New
York, 1973). {[}X6{]} K. Q. Brown, ``Voronoi diagrams from convex
hulls,'' Inf. Process. Lett. 9, 223--228 (1979). {[}X7{]} H.
Edelsbrunner, \emph{Geometry and Topology for Mesh Generation}
(Cambridge University Press, 2001). {[}X8{]} J. C. Maxwell, ``On the
calculation of the equilibrium and stiffness of frames,'' Phil. Mag. 27,
294--299 (1864). {[}X9{]} C. R. Calladine, ``Buckminster Fuller's
`tensegrity' structures and Clerk Maxwell's rules for the construction
of stiff frames,'' Int. J. Solids Struct. 14, 161--172 (1978). {[}X10{]}
G. Laman, ``On graphs and rigidity of plane skeletal structures,'' J.
Eng. Math. 4, 331--340 (1970). {[}X11{]} J. C. Phillips, ``Topology of
covalent non-crystalline solids I: Short-range order in chalcogenide
alloys,'' J. Non-Cryst. Solids 34, 153--181 (1979). {[}X12{]} M. F.
Thorpe, ``Continuous deformations in random networks,'' J. Non-Cryst.
Solids 57, 355--370 (1983). {[}X13{]} F. C. Frank, ``Supercooling of
liquids,'' Proc. R. Soc. Lond. A 215, 43--46 (1952). {[}X14{]} J.-F.
Sadoc, R. Mosseri, \emph{Geometrical Frustration} (Cambridge University
Press, 1999). {[}X15{]} D. R. Nelson, ``Order, frustration, and defects
in liquids and glasses,'' Phys. Rev.~B 28, 5515--5535 (1983). {[}X16{]}
M. Avrami, ``Kinetics of phase change. I. General theory,'' J. Chem.
Phys. 7, 1103--1112 (1939). {[}X17{]} A. H. Guth, ``Inflationary
universe: A possible solution to the horizon and flatness problems,''
Phys. Rev.~D 23, 347--356 (1981). {[}X18{]} A. H. Guth, E. J. Weinberg,
``Could the universe have recovered from a slow first-order phase
transition?,'' Nucl. Phys. B 212, 321--364 (1983). {[}X19{]} A. D.
Linde, ``A new inflationary universe scenario: A possible solution of
the horizon, flatness, homogeneity, isotropy and primordial monopole
problems,'' Phys. Lett. B 108, 389--393 (1982). {[}X20{]} T. W. B.
Kibble, ``Topology of cosmic domains and strings,'' J. Phys. A 9,
1387--1398 (1976). {[}X21{]} W. H. Zurek, ``Cosmological experiments in
superfluid helium?,'' Nature 317, 505--508 (1985). {[}X22{]} A. del
Campo, W. H. Zurek, ``Universality of phase transition dynamics:
Topological defects from symmetry breaking,'' Int. J. Mod. Phys. A 29,
1430018 (2014). {[}X23{]} T. Regge, ``General relativity without
coordinates,'' Nuovo Cimento 19, 558--571 (1961). {[}X24{]} J. Ambjørn,
J. Jurkiewicz, R. Loll, ``Emergence of a 4D world from causal quantum
gravity,'' Phys. Rev.~Lett. 93, 131301 (2004).

\begin{center}\rule{0.5\linewidth}{0.5pt}\end{center}

\subsection{Appendix A　Reproducing the
figures}\label{appendix-a-reproducing-the-figures}

Figures 1--9 were generated by
\texttt{make\_\allowbreak{}figures\_\allowbreak{}quaternion\_\allowbreak{}crystallization\_\allowbreak{}20261021.py} in the
same folder as this paper. Run
\texttt{python\ make\_\allowbreak{}figures\_\allowbreak{}quaternion\_\allowbreak{}crystallization\_\allowbreak{}20261021.py\ out};
dependencies are numpy and scipy (ConvexHull, Delaunay), and cairosvg
only for PNG conversion. The random seeds are 20261021 for Figure 5 (30
points) and 7 for Figure 6 (13 points). When run, Figure 5 checks and
prints the agreement between the projections of the faces of the convex
hull and the Delaunay triangulation of the projected points (agreement
on the run of 2026-10-21, 53 triangles). The closure model of Figure 8
(r\_n = 1 + 2(1 − V\_n/V\_max), V\_max = 10⁷⁸) is illustrative and not
derived. The English-label figures were generated by
\texttt{make\_\allowbreak{}figures\_\allowbreak{}quaternion\_\allowbreak{}crystallization\_\allowbreak{}20261021\_\allowbreak{}v1.1.py}
in the same folder with the argument \texttt{en} (\texttt{図表/\allowbreak{}en/\allowbreak{}}).
This program switches only the text of the figures between Japanese and
English, and its Japanese output is byte-for-byte identical to that of
the program above (the reproduction record is \texttt{図表/\allowbreak{}README.md}).

\subsection{Appendix B　Dialogue
record}\label{appendix-b-dialogue-record}

The content of this paper arose from a dialogue between the author and
Claude on 2026-10-21. The insights obtained in the dialogue, and the
points corrected or withdrawn, are summarized in
\texttt{追記\_\allowbreak{}対話での気づき\_\allowbreak{}研究動機と分解型八元数\_\allowbreak{}20261021.md} (in
Japanese) in the same folder.

\subsection{Version history}\label{version-history}

\begin{itemize}
\tightlist
\item
  v1 (2026-10-21): First draft (pre-publication draft; not published on
  Zenodo).
\item
  v1.1 (2026-10-21): The external reference Planck 2018 (v1's
  {[}X22{]}), which was not cited in the text, was removed from the
  references. The review del Campo--Zurek (2014) was added as {[}X22{]}
  as the basis for ``has been verified experimentally'' in §10.4, and
  Regge and Ambjørn et al.~(v1's {[}X23{]} and {[}X24{]}) were kept with
  the same numbers. The titles of the self-citations {[}S1{]}--{[}S6{]}
  in the Japanese version were made the same Japanese titles as in the
  Japanese version of the previous paper {[}S6{]} ({[}S4{]} with its
  full title), and the Concept DOIs were confirmed against the published
  Zenodo records. The bibliographic data of the external references were
  checked against the originals (Crossref, Zenodo); {[}X2{]} was changed
  to its Zenodo Concept DOI and the title of the latest version, and the
  description of {[}X2{]} in §12.1 was corrected, following the
  original, from ``free-energy minimization'' to ``energy
  minimization.'' The claims of the text are unchanged. This is the
  first version published on Zenodo (Version DOI
  10.5281/zenodo.23523238, Concept DOI 10.5281/zenodo.23523237).
\end{itemize}

\end{document}
